\documentclass{article}
\usepackage[utf8]{inputenc}
\usepackage[T1]{fontenc}
\usepackage{geometry}
\usepackage{graphicx}
\usepackage{float}
\usepackage{hyperref}
\usepackage[english,indonesian]{babel}

\title{Optimasi Pemrosesan Template LaTeX Menggunakan Pendekatan Berbasis AI}
\author{Academic Support Agent}
\date{\today}

\begin{document}
\maketitle

\begin{abstract}

\end{abstract}

\chapter{PENDAHULUAN}


\chapter{TINJAUAN PUSTAKA}


% Original content (commented out):
% \documentclass[a4paper, 12pt]{article}

\usepackage[a4paper,left=2cm, right=2cm,top=2cm, bottom=2cm]{geometry}
\usepackage{tikz}
\usepackage{tikzpagenodes} % nodes for tikz
\usepackage{lipsum}  
\usepackage{graphicx}
\usepackage{xcolor}
\usepackage{fontspec} % fonts
\usepackage{hyperref}%
\usepackage{vwcol} % multi column
\usepackage{float} % table position
\usepackage{booktabs}

\usepackage[
backend=biber,
style=numeric,
sorting=nyt,
natbib=true,
maxbibnames=9, 
maxcitenames=2
]{biblatex} % include biblatex
\addbibresource{bibliography.bib} % include .bib file
\DeclareNameAlias{author}{last-first} % lastname first

\setlength\parindent{0pt} % set indent 0
\usetikzlibrary{positioning, calc} % setup tikz

\begin{document}
\input{sources/frontpage/frontpage} % input frontpage
\input{sources/test}
\newpage
\printbibliography 
\end{document}



\end{document}